\chapter{Espaces complets~: Baire, Ascoli, Stone--Weierstrass}
\label{ch:b3:complete}

La complétude --- toute suite de Cauchy converge --- est la
propriété qui permet à l'analyse de \emph{produire} des
objets~: points fixes de contractions, sommes de séries,
solutions d'équations obtenues comme limites. Ce chapitre
assemble les trois grandes machines d'existence de la théorie
métrique. Le \emph{théorème de Baire} montre qu'un espace
\omterm{def:b3:complete:complete}{complet} ne peut être réunion dénombrable de morceaux
négligeables, et fait apparaître des objets (fonctions \omterm{def:b3:topology:continuity}{continues}
nulle part dérivables~!) par un pur raisonnement de cardinalité.
\emph{Arzelà--Ascoli} identifie les parties compactes de
$\mathcal C(K)$ et est le cheval de bataille de la compacité en
analyse --- le problème de fin de semaine s'en sert pour prouver
le théorème d'existence de Peano pour les équations
différentielles. \emph{Stone--Weierstrass} montre que les
polynômes, et bien d'autres, sont \omterm{def:b3:topology:interior}{denses} dans $\mathcal C(K)$~:
l'approximation devient une vérification \omterm{def:b3:galois:algebraic}{algébrique}. En chemin
on construit les \omterm{thm:b3:complete:completion}{complétés} et on prouve le théorème
d'extension des applications uniformément \omterm{def:b3:topology:continuity}{continues}, le pain
quotidien des~\cref{ch:b3:lp,ch:b3:hilbert,ch:b3:fouriertransform}.

\section{Espaces complets, complétés, extensions}

\begin{definition}\label{def:b3:complete:complete}
Un espace métrique est \emph{complet}\index{espace métrique
complet} si toute suite de Cauchy converge (deuxième année~:
$\R^n$ est complet~; $\mathcal C(\intcc01)$ muni de
$\norm\cdot_\infty$ est complet). Une partie fermée d'un
espace complet est complète~; une partie complète d'un espace
métrique quelconque est fermée.
\end{definition}

\begin{proof}
Pour les deux énoncés~: une suite de Cauchy du fermé $F$
converge dans $X$, et sa limite, adhérente à $F$, est dans
$F$~; une suite convergente-dans-$X$ d'une partie complète $A$
est de Cauchy, donc converge dans $A$, et les limites sont
uniques.
\end{proof}

\begin{theorem}[Extension des applications uniformément continues]
\label{thm:b3:complete:extension}
Soit $D \subseteq X$ \omterm{def:b3:topology:interior}{dense}, $Y$ \omterm{def:b3:complete:complete}{complet}, et $f \colon D \to
Y$ uniformément \omterm{def:b3:topology:continuity}{continue}. Alors $f$ s'étend \emph{uniquement}
en une $\bar f \colon X \to Y$ \omterm{def:b3:topology:continuity}{continue}, et $\bar f$ est
uniformément \omterm{def:b3:topology:continuity}{continue}.\index{théorème d'extension (continuité
uniforme)}
\end{theorem}

\begin{proof}
Unicité~: deux extensions \omterm{def:b3:topology:continuity}{continues} coïncident sur le \omterm{def:b3:topology:interior}{dense}
$D$, donc partout (l'ensemble d'accord $\{g = h\}$ est fermé~:
image réciproque de la diagonale fermée sous $x\mapsto(g(x),
h(x))$).
Existence~: pour $x \in X$ choisir $d_n \to x$, $d_n \in D$.
La suite $(f(d_n))$ est de Cauchy~: étant donné $\varepsilon$,
la continuité uniforme fournit $\delta$ avec $d(u,v) < \delta
\Rightarrow d(f(u), f(v)) < \varepsilon$, et $(d_n)$ est de
Cauchy. Définir $\bar f(x) = \lim f(d_n)$~; la limite ne
dépend pas de la suite choisie (entrelacer deux d'entre elles).
$\bar f$ prolonge $f$ (suites constantes) et hérite du \omterm{def:b3:modules:module}{module}
de continuité~: si $d(x, x') < \delta$, en approximant les
deux par des points de $D$ à distance $< \frac{\delta -
d(x,x')}2$ on obtient $d(\bar f(x), \bar f(x')) \leq
\varepsilon$ à la limite --- $\bar f$ est uniformément
\omterm{def:b3:topology:continuity}{continue}.
\end{proof}

\begin{theorem}[Complété]\label{thm:b3:complete:completion}
Tout espace métrique $X$ s'immerge isométriquement comme
\omterm{def:b3:topology:interior}{partie dense} d'un \omterm{def:b3:complete:complete}{espace métrique complet} $\hat X$, unique à
isométrie près fixant $X$ point par point~: son
\emph{complété}\index{complété d'un espace métrique}.
\end{theorem}

\begin{proof}
\emph{Existence.} Soit $\mathcal X$ l'ensemble des suites de
Cauchy de $X$, muni de la pseudo-distance
\[
D\bigl((x_n), (y_n)\bigr) = \lim_n d(x_n, y_n),
\]
la limite existant car $\abs{d(x_n, y_n) - d(x_m, y_m)}
\leq d(x_n, x_m) + d(y_n, y_m)$ rend la suite réelle de
Cauchy. Soit $\hat X = \mathcal X/{\sim}$, en identifiant les
suites à $D$-distance $0$~; $D$ descend en une distance.
Immerger $X$ par les suites constantes~: une isométrie, d'image
\omterm{def:b3:topology:interior}{dense} (une suite de Cauchy est $D$-approximée par les
constantes construites sur ses propres termes~:
$D\bigl((x_n), (x_k)_{\rm const}\bigr) = \lim_n d(x_n, x_k)
\to 0$ quand $k \to \infty$ par Cauchy). Complétude de $\hat
X$~: soit $(\xi^k)$ de Cauchy dans $\hat X$~; par densité
choisir $x_k \in X$ avec $D(\xi^k, x_k) \leq 2^{-k}$~; alors
$(x_k)$ est de Cauchy dans $X$ (inégalité triangulaire via les
$\xi$), définit un point $\xi \in \hat X$, et $D(\xi^k, \xi)
\leq 2^{-k} + D(x_k, \xi) \to 0$ (la distance de la constante
$x_k$ à la classe de $(x_j)_j$ est $\lim_j d(x_k, x_j)$, petite
pour $k$ grand).

\emph{Unicité}~: deux \omterm{thm:b3:complete:completion}{complétés} $\hat X_1, \hat X_2$
contiennent $X$ densément~; l'identité de $X$, une isométrie,
est uniformément \omterm{def:b3:topology:continuity}{continue}, donc s'étend en $\hat X_1 \to \hat
X_2$ (\cref{thm:b3:complete:extension}), encore une isométrie
sur un ensemble \omterm{def:b3:topology:interior}{dense} donc partout~; symétriquement dans
l'autre sens, et les composées fixent le \omterm{def:b3:topology:interior}{dense} $X$~: ce sont
les identités.
\end{proof}

\begin{theorem}[Point fixe de Banach]\label{thm:b3:complete:banach}
Soit $X$ \omterm{def:b3:complete:complete}{complet}, non vide, et $f \colon X \to X$ une
\emph{contraction}~: $d(f(x), f(y)) \leq k\,d(x,y)$ avec $k <
1$. Alors $f$ a un unique point fixe $x^*$, et toute \omterm{def:b3:groups:action}{orbite}
converge vers lui, avec le taux explicite $d(x_n, x^*) \leq
\frac{k^n}{1-k}\,d(x_1, x_0)$.\index{théorème du point fixe de
Banach}
\end{theorem}

\begin{proof}
(La deuxième année l'a prouvé~; on réenregistre l'argument en
deux lignes pour l'autonomie.) L'\omterm{def:b3:groups:action}{orbite} $x_{n+1} = f(x_n)$
vérifie $d(x_{n+1}, x_n) \leq k^nd(x_1, x_0)$, donc est de
Cauchy (série géométrique)~; sa limite $x^*$ est fixe
(continuité de $f$), unique car deux points fixes satisfont $d
\leq k\,d$. Le taux~: sommer la queue géométrique.
\end{proof}

\begin{example}[Perturber l'identité]
\label{ex:b3:complete:perturbident}
Soit $g \colon \R^d \to \R^d$ $k$-lipschitzienne avec $k < 1$.
Alors $\varphi = \mathrm{id} + g$ est un \emph{\omterm{def:b3:topology:continuity}{homéomorphisme}
de $\R^d$ sur $\R^d$}. Injectivité, avec un \omterm{def:b3:modules:module}{module} quantitatif~:
\[
\norm{\varphi(x) - \varphi(y)} \geq \norm{x - y} -
\norm{g(x) - g(y)} \geq (1 - k)\norm{x - y} .
\]
La surjectivité est le théorème du point fixe~: résoudre
$\varphi(x) = y$ signifie $x = y - g(x)$, et $x \mapsto y -
g(x)$ est une $k$-contraction du \omterm{def:b3:complete:complete}{complet} $\R^d$ --- une
solution unique $x = \psi(y)$ existe pour tout $y$.
L'inégalité affichée rend l'inverse $\psi$ lipschitzien de
constante $\frac1{1-k}$~: un \omterm{def:b3:topology:continuity}{homéomorphisme}, avec bornes
explicites sur les deux \omterm{def:b3:modules:module}{modules}. Cet énoncé d'apparence
innocente est le moteur du théorème d'inversion locale
(\cref{ch:b3:submanifolds})~: près d'un point où $Df$ est
inversible, $f$ \emph{est} une application linéaire inversible
plus une petite perturbation lipschitzienne, et l'exemple
d'aujourd'hui fait le reste. Il quantifie aussi la robustesse
numérique~: un système perturbé de moins que la marge de
l'inverse reste solvable, la solution se déplaçant d'au plus
$\frac{1}{1-k}$ fois la perturbation.
\end{example}

\section{Le théorème de Baire}

\begin{theorem}[Baire]\label{thm:b3:complete:baire}
Dans un \omterm{def:b3:complete:complete}{espace métrique complet}, une intersection dénombrable
d'\omterm{def:b3:topology:topology}{ouverts} \omterm{def:b3:topology:interior}{denses} est \omterm{def:b3:topology:interior}{dense}. De façon équivalente~: si $X =
\bigcup_{n}F_n$ avec chaque $F_n$ fermé, alors quelque $F_n$ a
un intérieur non vide.\index{théorème de Baire}
\end{theorem}

\begin{proof}
Soit $(U_n)$ des \omterm{def:b3:topology:topology}{ouverts} \omterm{def:b3:topology:interior}{denses} et $B_0 = B(x_0, r_0)$ une
boule ouverte quelconque~; on trouve un point de $\bigcap U_n$
dans $B_0$. Inductivement~: $U_{n}$ étant \omterm{def:b3:topology:interior}{dense} et \omterm{def:b3:topology:topology}{ouvert}, il
rencontre la boule ouverte $B_{n-1}$ en un \omterm{def:b3:topology:topology}{ouvert}, qui contient
une boule fermée $\bar B(x_n, r_n)$ avec $0 < r_n \leq
r_{n-1}/2$ et $\bar B(x_n, r_n) \subseteq B_{n-1}\cap U_n$.
Les \omterm{ex:b3:groups:actions}{centres} forment une suite de Cauchy ($x_m \in B_n$ pour $m
\geq n$, rayons $\to 0$)~; la limite $x$ est dans chaque $\bar
B(x_n, r_n)$ (fermeture), donc dans chaque $U_n$ et dans $B_0$.
Pour la seconde forme~: si aucun $F_n$ n'a d'intérieur, les
$U_n = X \setminus F_n$ sont \omterm{def:b3:topology:topology}{ouverts} et \omterm{def:b3:topology:interior}{denses}, et un point de
$\bigcap U_n$ échappe à $\bigcup F_n = X$~: absurde.
\end{proof}

\begin{remark}\label{rem:b3:complete:meagre}
Vocabulaire~: un ensemble est \emph{nulle part \omterm{def:b3:topology:interior}{dense}} si son
\omterm{def:b3:topology:interior}{adhérence} a un intérieur vide, \emph{maigre}\index{ensemble
maigre} (première catégorie) s'il est réunion dénombrable
d'ensembles nulle part \omterm{def:b3:topology:interior}{denses}. Baire~: \emph{un \omterm{def:b3:complete:complete}{espace métrique
complet} n'est pas \omterm{rem:b3:complete:meagre}{maigre} en lui-même}, et le complémentaire
d'un \omterm{rem:b3:complete:meagre}{ensemble maigre} est \omterm{def:b3:topology:interior}{dense}. «~\omterm{rem:b3:complete:meagre}{Maigre}~» est une notion de
petitesse orthogonale à la mesure (le~\cref{ch:b3:measure}
produira des \omterm{rem:b3:complete:meagre}{ensembles maigres} de mesure pleine), et les
arguments de Baire prouvent l'\emph{existence par abondance}~:
pour exhiber un objet sans la propriété P, montrer que les
objets-P forment un \omterm{rem:b3:complete:meagre}{ensemble maigre}.
\end{remark}

\begin{corollary}\label{cor:b3:complete:bairecors}
(a) $\R$ est non dénombrable. (b) $\Q$ n'est pas une
intersection dénombrable d'\omterm{def:b3:topology:topology}{ouverts} de $\R$, et un \omterm{def:b3:complete:complete}{espace
métrique complet} non vide sans points isolés est non
dénombrable.
\end{corollary}

\begin{proof}
(a) $\R = \bigcup_{x}\{x\}$ sur un ensemble dénombrable
rendrait quelque singleton d'intérieur non vide. (b) Si $\Q =
\bigcap_n V_n$ avec $V_n$ \omterm{def:b3:topology:topology}{ouverts} (nécessairement \omterm{def:b3:topology:interior}{denses}, car
$\supseteq \Q$), alors les $V_n$ et les complémentaires
$\R\setminus\{q\}$, $q \in \Q$, forment une famille dénombrable
d'\omterm{def:b3:topology:topology}{ouverts} \omterm{def:b3:topology:interior}{denses} d'intersection vide --- contredisant Baire.
Si $X$ est \omterm{def:b3:complete:complete}{complet} sans points isolés et dénombrable, $X =
\bigcup_{x \in X}\{x\}$ l'exhibe comme réunion dénombrable de
fermés d'intérieur vide (pas de points isolés)~: Baire encore.
\end{proof}

\begin{theorem}[Les monstres de Weierstrass existent]
\label{thm:b3:complete:nowherediff}
Il existe des fonctions \omterm{def:b3:topology:continuity}{continues} sur $\intcc01$ dérivables en
aucun point. Mieux~: l'ensemble des $f \in \mathcal
C(\intcc01)$ ayant une dérivée (finie) en au moins un point est
\omterm{rem:b3:complete:meagre}{maigre} dans $(\mathcal C(\intcc01), \norm\cdot_\infty)$.%
\index{fonction nulle part dérivable}
\end{theorem}

\begin{proof}
Pour $n \geq 1$ soit
\[
F_n = \Bigl\{f : \exists x \in \intcc01,\
\forall h \neq 0 \text{ avec } x + h \in \intcc01,\
\abs{f(x+h) - f(x)} \leq n\abs h \Bigr\}.
\]
Si $f$ est dérivable en $x$, alors $f \in F_n$ pour quelque
$n$~: le quotient $\abs{f(x+h)-f(x)}/\abs h$ est borné pour
$\abs h \leq \delta$ (dérivabilité~: il tend vers
$\abs{f'(x)}$) et borné par $2\norm f_\infty/\delta$ pour
$\abs h \geq \delta$. Donc $\bigcup F_n$ contient toutes les
fonctions quelque part dérivables, et il suffit de montrer
chaque $F_n$ fermé d'intérieur vide.

\emph{Fermé}~: soit $f_k \to f$ uniformément, $f_k \in F_n$
avec témoins $x_k \to x$ (compacité, après extraction). Pour
$h$ avec $x + h \in \intcc01$~: choisir $h_k \to h$ avec $x_k
+ h_k \in \intcc01$ (p.\,ex.\ $h_k = h + x - x_k$ tronqué)~;
alors $\abs{f(x + h) - f(x)} = \lim \abs{f_k(x_k + h_k) -
f_k(x_k)} \leq \lim n\abs{h_k} = n\abs h$, en utilisant la
convergence uniforme et la continuité de $f$ aux points
pertinents~: $f \in F_n$.

\emph{Intérieur vide}~: étant donné $f \in F_n$ et
$\varepsilon > 0$, on trouve $g$ avec $\norm{g - f}_\infty \leq
\varepsilon$ et $g \notin F_n$. D'abord approximer $f$ à
$\varepsilon/2$ près par une $\varphi$ affine par morceaux
(continuité uniforme~: interpoler sur une grille fine), de
pentes bornées par quelque $M$. Ajouter une petite dent de
scie~: $g = \varphi + \frac\varepsilon2\,s_N$, où $s_N(x)$ est
le zigzag $\frac1N$-périodique d'amplitude $1$ et de pente
$\pm 2N$. En tout $x$, d'un côté il y a $h$ arbitrairement
petit avec la dent de scie contribuant une pente $\pm
2N\cdot\frac\varepsilon2 = \pm\varepsilon N$ sur $[x, x+h]$~:
le quotient différentiel de $g$ dépasse $\varepsilon N - M > n$
pour $N$ grand. Donc $g \notin F_n$, à toute distance uniforme
$\leq \varepsilon$ de $f$. Conclusion~: $\bigcup F_n$ est
\omterm{rem:b3:complete:meagre}{maigre}~; par Baire son complémentaire --- formé de \omterm{thm:b3:complete:nowherediff}{fonctions
nulle part dérivables} --- est \omterm{def:b3:topology:interior}{dense} dans $\mathcal
C(\intcc01)$~: de telles fonctions existent \emph{en
abondance}.
\end{proof}

\section{Arzelà--Ascoli}

Tout au long, $K$ est un espace métrique \emph{compact} et
$\mathcal C(K) = \mathcal C(K, \R^d)$, avec $\norm f_\infty =
\sup_K \norm{f(x)}$~: un espace \omterm{def:b3:complete:complete}{complet} (les limites uniformes
de \omterm{def:b3:topology:continuity}{continues} sont \omterm{def:b3:topology:continuity}{continues} --- deuxième année).

\begin{definition}\label{def:b3:complete:equicontinuous}
Une famille $\mathcal F \subseteq \mathcal C(K)$ est
\emph{équicontinue}\index{équicontinuité} si pour tout
$\varepsilon > 0$ il existe $\delta > 0$ tel que
\[
d(x, y) < \delta \implies \norm{f(x) - f(y)} < \varepsilon
\quad\text{pour \emph{toutes} les } f \in \mathcal F
\]
(un seul $\delta$ pour toute la famille --- p.\,ex.\ toute
famille avec une constante de Lipschitz commune, ou un \omterm{def:b3:modules:module}{module}
de Hölder commun), et \emph{bornée ponctuellement} si
$\sup_{f}\norm{f(x)} < \infty$ pour chaque $x$.
\end{definition}

\begin{theorem}[Arzelà--Ascoli]\label{thm:b3:complete:ascoli}
Une partie $\mathcal F \subseteq \mathcal C(K)$ est
relativement compacte (d'\omterm{def:b3:topology:interior}{adhérence} compacte) ssi elle est
\omterm{def:b3:complete:equicontinuous}{équicontinue} et bornée ponctuellement. En particulier, toute
suite \omterm{def:b3:complete:equicontinuous}{équicontinue} et bornée ponctuellement a une sous-suite
uniformément convergente.\index{théorème d'Arzelà--Ascoli}
\end{theorem}

\begin{proof}
($\Leftarrow$) Soit $(f_k)$ une suite de $\mathcal F$. $K$
métrique compact est \emph{\omterm{def:b3:galois:separable}{séparable}}~: pour chaque $n$, un
nombre fini de boules de rayon $\frac1n$ recouvrent $K$
(\cref{thm:b3:topology:metriccompact})~; leurs \omterm{ex:b3:groups:actions}{centres} forment
un ensemble \omterm{def:b3:topology:interior}{dense} dénombrable $D = \{x_1, x_2, \dots\}$. Par
bornitude ponctuelle et Bolzano--Weierstrass, extraire
successivement des sous-suites convergeant en $x_1$, puis aussi
en $x_2$, etc., et prendre la sous-suite \emph{diagonale}
$(g_j)$~: elle converge en tout point de $D$. L'\omterm{def:b3:complete:equicontinuous}{équicontinuité}
promouvoit cela en Cauchy uniforme~: étant donné
$\varepsilon$, prendre $\delta$ comme dans la définition,
recouvrir $K$ par un nombre fini de boules $B(x_i, \delta)$
avec $x_i \in D$ ($i \leq m$), et choisir $J$ assez grand pour
que $\norm{g_j(x_i) - g_{j'}(x_i)} < \varepsilon$ pour $j, j'
\geq J$, $i \leq m$. Pour $x \in B(x_i, \delta)$ arbitraire~:
\[
\norm{g_j(x) - g_{j'}(x)}
\leq \norm{g_j(x) - g_j(x_i)} + \norm{g_j(x_i) - g_{j'}(x_i)}
+ \norm{g_{j'}(x_i) - g_{j'}(x)} < 3\varepsilon .
\]
Donc $(g_j)$ est uniformément de Cauchy, et converge dans le
\omterm{def:b3:complete:complete}{complet} $\mathcal C(K)$. Ainsi toute suite de $\mathcal F$ a
une sous-suite convergente~: $\bar{\mathcal F}$ est
(séquentiellement, donc par
\cref{thm:b3:topology:metriccompact}) compacte.

($\Rightarrow$) Si $\bar{\mathcal F}$ est compacte~: la
bornitude ponctuelle est claire (l'évaluation est \omterm{def:b3:topology:continuity}{continue}).
Pour l'\omterm{def:b3:complete:equicontinuous}{équicontinuité}, recouvrir $\mathcal F$ par un nombre
fini de boules $B(f_i, \varepsilon)$ de $\mathcal C(K)$~;
chaque $f_i$ est uniformément \omterm{def:b3:topology:continuity}{continue} (Heine,
\cref{cor:b3:topology:heineborel}), donnant un $\delta$ commun
pour $i \leq m$~; alors pour $f \in B(f_i, \varepsilon)$ et
$d(x,y) < \delta$~: $\norm{f(x)-f(y)} \leq 2\varepsilon +
\norm{f_i(x)-f_i(y)} < 3\varepsilon$.
\end{proof}

\begin{example}\label{ex:b3:complete:ascoliexamples}
La boule unité fermée de $\mathcal C(\intcc01)$ n'est
\emph{pas} compacte ($f_n(x) = x^n$ n'a pas de sous-suite
uniformément convergente~: la limite ponctuelle est
discontinue), et en effet $(x^n)$ n'est pas \omterm{def:b3:complete:equicontinuous}{équicontinue} en
$1$. En revanche $\{f : \norm f_\infty \leq 1,\
\operatorname{Lip}(f) \leq 1\}$ est compacte~: bornée et
$1$-lipschitz-équicontinue, et fermée. Ascoli explique
\emph{pourquoi} la compacité échoue en dimension infinie
(Riesz, deuxième année) et quoi ajouter pour la restaurer~: un
\omterm{def:b3:modules:module}{module} de continuité uniforme.
\end{example}

\section{Stone--Weierstrass}

\begin{lemma}[Dini]\label{lem:b3:complete:dini}
Soit $K$ compact et $(f_n)$ une suite \emph{monotone} de
fonctions réelles \omterm{def:b3:topology:continuity}{continues} convergeant \emph{ponctuellement}
vers une $f$ \omterm{def:b3:topology:continuity}{continue}. Alors la convergence est
uniforme.\index{théorème de Dini}
\end{lemma}

\begin{proof}
Disons $f_n \uparrow f$~; soit $g_n = f - f_n \downarrow 0$,
\omterm{def:b3:topology:continuity}{continues}. Étant donné $\varepsilon$, les \omterm{def:b3:topology:topology}{ouverts} $U_n =
\{g_n < \varepsilon\}$ croissent et recouvrent $K$
(convergence ponctuelle)~; extraire un sous-recouvrement
fini~: $K = U_{n_0}$ pour quelque $n_0$ (famille croissante),
c'est-à-dire $0 \leq g_n < \varepsilon$ partout pour $n \geq
n_0$.
\end{proof}

\begin{lemma}\label{lem:b3:complete:sqrt}
Il existe une suite de \emph{polynômes} $u_n$ avec $u_n(t) \to
\sqrt t$ uniformément sur $\intcc01$.
\end{lemma}

\begin{proof}
Définir $u_0 = 0$, $u_{n+1}(t) = u_n(t) + \frac12\bigl(t -
u_n(t)^2\bigr)$~: des polynômes. Par induction $0 \leq u_n(t)
\leq \sqrt t$ sur $\intcc01$~: en l'accordant pour $n$,
\[
\sqrt t - u_{n+1}(t) = \bigl(\sqrt t - u_n(t)\bigr)
\Bigl(1 - \tfrac12\bigl(\sqrt t + u_n(t)\bigr)\Bigr) \geq 0,
\]
puisque $\sqrt t + u_n \leq 2$~; et $u_{n+1} \geq u_n \geq 0$.
Donc $(u_n(t))$ est croissante, bornée par $\sqrt t$~: elle
converge ponctuellement, et la limite $\ell(t)$ satisfait $\ell
= \ell + \frac12(t - \ell^2)$~: $\ell(t) = \sqrt t$, \omterm{def:b3:topology:continuity}{continue}.
Dini (\cref{lem:b3:complete:dini}) promeut en uniforme.
\end{proof}

\begin{theorem}[Stone--Weierstrass, version réelle]
\label{thm:b3:complete:stoneweierstrass}
Soit $K$ un \omterm{def:b3:topology:compact}{espace compact} (\omterm{def:b3:topology:hausdorff}{séparé}) et $\mathcal A \subseteq
\mathcal C(K, \R)$ une \emph{sous-algèbre} (stable par sommes,
produits, multiplications scalaires) qui \emph{contient les
constantes} et \emph{sépare les points} (pour $x \neq y$,
quelque $f \in \mathcal A$ a $f(x) \neq f(y)$). Alors $\mathcal
A$ est \omterm{def:b3:topology:interior}{dense} dans $(\mathcal C(K,\R), \norm\cdot_\infty)$.%
\index{théorème de Stone--Weierstrass}
\end{theorem}

\begin{proof}
Soit $\bar{\mathcal A}$ l'\omterm{def:b3:topology:interior}{adhérence}, encore une algèbre (les
produits de limites uniformes sur des ensembles bornés
convergent).

\emph{Étape 1~: $\bar{\mathcal A}$ est un treillis}, c'est-à-dire
stable par $\max$ et $\min$. Comme $\max(f,g) = \frac{f + g +
\abs{f-g}}2$ et $\min$ de même, il suffit que $f \in
\bar{\mathcal A} \Rightarrow \abs f \in \bar{\mathcal A}$~: avec
$M = \norm f_\infty > 0$, $\abs f = M\sqrt{(f/M)^2}$, et
\cref{lem:b3:complete:sqrt} donne des polynômes $u_n$ avec
$u_n\bigl((f/M)^2\bigr) \to \abs f/M$ uniformément~; les
polynômes en des membres de l'algèbre (avec terme constant~: les
constantes y sont) restent dans $\bar{\mathcal A}$.

\emph{Étape 2~: interpolation en deux points.} Pour $x \neq y$
et $a, b \in \R$, quelque $g \in \mathcal A$ a $g(x) = a$,
$g(y) = b$~: prendre $h$ séparant $x, y$ et poser $g = a + (b -
a)\frac{h - h(x)}{h(y) - h(x)}$.

\emph{Étape 3.} Soit $f \in \mathcal C(K)$, $\varepsilon > 0$.
Pour chaque paire $x, y$ choisir $g_{x,y} \in \mathcal A$ avec
$g_{x,y}(x) = f(x)$, $g_{x,y}(y) = f(y)$ (étape 2~; pour $x =
y$ prendre la constante $g_{x,x} = f(x)$).
Fixer $x$~: pour chaque $y$, l'\omterm{def:b3:topology:topology}{ouvert} $V_y = \{g_{x,y} < f +
\varepsilon\}$ contient $y$~; la compacité extrait $y_1, \dots,
y_m$ avec $K = \bigcup V_{y_j}$, et $h_x = \min_j g_{x, y_j}
\in \bar{\mathcal A}$ (étape 1) vérifie $h_x < f +
\varepsilon$ partout, $h_x(x) = f(x)$. Faire varier $x$~: $W_x
= \{h_x > f - \varepsilon\}$ est \omterm{def:b3:topology:topology}{ouvert}, contient $x$~; extraire
$x_1, \dots, x_l$ recouvrant $K$, et $h = \max_i h_{x_i} \in
\bar{\mathcal A}$ vérifie $f - \varepsilon < h < f +
\varepsilon$~: $\norm{f - h}_\infty \leq \varepsilon$. Donc $f
\in \bar{\mathcal A}$.
\end{proof}

\begin{corollary}\label{cor:b3:complete:weierstrass}
(a) (\emph{Weierstrass}) Les polynômes sont \omterm{def:b3:topology:interior}{denses} dans
$\mathcal C(\intcc ab, \R)$~; les polynômes en $n$ variables
sont \omterm{def:b3:topology:interior}{denses} dans $\mathcal C(K, \R)$ pour $K \subseteq \R^n$
compact.
(b) (\emph{Version complexe}) Si $\mathcal A \subseteq \mathcal
C(K, \C)$ est une sous-algèbre contenant les constantes,
séparant les points, et \emph{stable par conjugaison}, elle est
\omterm{def:b3:topology:interior}{dense}.
(c) (\emph{Version trigonométrique}) Les polynômes
trigonométriques $\sum_{\abs n \leq N}c_n\eu^{\iu n t}$ sont
\omterm{def:b3:topology:interior}{denses} dans l'espace des fonctions \omterm{def:b3:topology:continuity}{continues} $2\pi$-périodiques
muni de $\norm\cdot_\infty$.\index{théorème d'approximation de
Weierstrass}
\end{corollary}

\begin{proof}
(a) Les polynômes forment une algèbre avec constantes~; les
fonctions coordonnées séparent les points de $\R^n$.
(b) Les parties réelle et imaginaire $\frac{f + \bar f}2$,
$\frac{f - \bar f}{2\iu}$ des membres de $\mathcal A$ forment
une algèbre réelle $\mathcal A_\R \subseteq \mathcal C(K, \R)$
avec constantes~; elle sépare les points ($f(x) \ne f(y)$ force
$\operatorname{Re}f$ ou $\operatorname{Im}f$ à séparer).
Appliquer le théorème réel et recombiner.
(c) Voir les fonctions $2\pi$-périodiques comme $\mathcal
C(S^1, \C)$ ($S^1 = \R/2\pi\Z$, compact~:
\cref{exo:b3:topology:5})~; l'algèbre engendrée par
$\eu^{\iu t}, \eu^{-\iu t}$ et les constantes est stable par
conjugaison et sépare les points du cercle ($\eu^{\iu t}$ est
injective dessus). Appliquer (b).
\end{proof}

\begin{remark}\label{rem:b3:complete:swremark}
La version trigonométrique répare, et généralise largement, le
trou laissé dans le chapitre Fourier de deuxième année~: la
densité des polynômes trigonométriques dans $(\mathcal
C(S^1), \norm\cdot_2)$ suit a fortiori ($\norm\cdot_2 \leq
\norm\cdot_\infty$ à la constante de normalisation près), ce
qui fera du système de Fourier une \emph{\omterm{def:b3:topology:basis}{base}} orthonormée
dans le~\cref{ch:b3:hilbert}, prouvant enfin Parseval en toute
généralité.
\end{remark}

\section{Exercices}

\begin{exercise}[$\star$]\label{exo:b3:complete:1}
(a) Montrer que $\mathcal C(\intcc01, \R)$ muni de
$\norm f_1 = \int_0^1\abs f$ n'est \emph{pas} \omterm{def:b3:complete:complete}{complet}~: les
fonctions $f_n$ égales à $0$ sur $\intcc0{\frac12}$, $1$ sur
$\intcc{\frac12 + \frac1n}1$, affines entre, sont
$\norm\cdot_1$-Cauchy sans limite \omterm{def:b3:topology:continuity}{continue}.
(b) Montrer qu'un espace normé dans lequel toute série
absolument convergente converge est \omterm{def:b3:complete:complete}{complet}. \emph{(Extraire
d'une suite de Cauchy une sous-suite avec $\norm{x_{n_{k+1}} -
x_{n_k}} \leq 2^{-k}$.)}
\end{exercise}

\begin{exercise}[$\star$]\label{exo:b3:complete:2}
En utilisant Baire~: (a) montrer qu'un espace normé \omterm{def:b3:complete:complete}{complet} n'a
pas de \omterm{def:b3:topology:basis}{base} (\omterm{def:b3:galois:algebraic}{algébrique}) dénombrable --- en déduire que
l'espace des polynômes n'est \omterm{def:b3:complete:complete}{complet} pour \emph{aucune}
norme~; (b) montrer que si une suite de $f_n \colon \R \to \R$
\omterm{def:b3:topology:continuity}{continues} converge ponctuellement vers $f$, l'ensemble des
points de continuité de $f$ est \omterm{def:b3:topology:interior}{dense}.
\emph{(Pour (b)~: admettre ou prouver que $\Omega_\delta =
\{x : \operatorname{osc}_x f < \delta\}$ est \omterm{def:b3:topology:topology}{ouvert}, et
montrer qu'il est \omterm{def:b3:topology:interior}{dense} en utilisant $F_{N} = \{x:
\abs{f_n(x) - f_m(x)} \leq \delta/3\ \forall n,m \geq N\}$,
fermés recouvrant $\R$~; travailler dans une boule fermée
arbitraire pour appliquer Baire.)}
\end{exercise}

\begin{exercise}[$\star\star$]\label{exo:b3:complete:3}
(a) Montrer que $f(x) = \frac12\bigl(x + \frac ax\bigr)$ ($a >
1$) est une contraction de $[\sqrt a, +\infty)$ et identifier
son point fixe --- méthode d'Héron. Estimer le nombre
d'itérations pour une précision $10^{-12}$ partant de $x_0 =
a$, pour $a = 2$.
(b) (Équation de Kepler) Pour $0 \leq e < 1$ et $m \in \R$,
montrer que $x = m + e\sin x$ a une unique solution, dépendant
continûment de $m$.
\end{exercise}

\begin{exercise}[$\star\star$]\label{exo:b3:complete:4}
(a) Soit $X$ \omterm{def:b3:complete:complete}{complet} et $f \colon X \to X$ telle qu'une
itérée $f^p$ est une contraction. Montrer que $f$ a un unique
point fixe. Application~: l'opérateur intégral $T$ sur
$\mathcal C(\intcc0a)$, $Tf(x) = \int_0^x f$, satisfait
$\norm{T^p} \leq a^p/p!$ --- en déduire que $u = g + \lambda
Tu$ est solvable pour tout $\lambda$.
(b) (Edelstein) Soit $K$ \emph{compact} et $f\colon K \to K$
avec $d(f(x), f(y)) < d(x, y)$ pour $x \neq y$. Montrer que
$f$ a un unique point fixe, mais que le taux de contraction
peut se perdre~: sur $X = [1, +\infty)$ (\omterm{def:b3:complete:complete}{complet}, non
compact), $f(x) = x + \frac1x$ n'a pas de point fixe malgré
des distances strictement décroissantes.
\end{exercise}

\begin{exercise}[$\star\star$]\label{exo:b3:complete:5}
(a) Deux \omterm{def:b3:topology:continuity}{applications continues} vers un \omterm{def:b3:topology:hausdorff}{espace séparé}
coïncidant sur une \omterm{def:b3:topology:interior}{partie dense} coïncident partout~; où cela
a-t-il été utilisé dans le chapitre~
(b) Soit $D \subseteq X$ \omterm{def:b3:topology:interior}{dense} et $f \colon D \to Y$ une
bijection isométrique sur une \omterm{def:b3:topology:interior}{partie dense} d'un $Y$ \omterm{def:b3:complete:complete}{complet},
avec $X$ \omterm{def:b3:complete:complete}{complet}. Montrer que $f$ s'étend en une bijection
isométrique $X \to Y$. En déduire à nouveau l'unicité des
\omterm{thm:b3:complete:completion}{complétés}.
\end{exercise}

\begin{exercise}[$\star\star$]\label{exo:b3:complete:6}
Lesquelles des familles suivantes sont \omterm{def:b3:complete:equicontinuous}{équicontinues}, bornées
ponctuellement, relativement compactes dans $\mathcal
C(\intcc01)$~
\[
\{x \mapsto \sin(nx)\}_{n},\qquad
\{x \mapsto x^n\}_n,
\]
\[
\Bigl\{f \in \mathcal C^1 : \norm f_\infty \leq 1,\
\norm{f'}_\infty \leq 1\Bigr\},\qquad
\{f : \operatorname{Lip}(f)\leq 1,\ f(0) = 0\}.
\]
Justifier chaque réponse par Ascoli ou une suite
contre-exemple.
\end{exercise}

\begin{exercise}[$\star\star\star$]\label{exo:b3:complete:7}
(Opérateurs intégraux compacts) Soit $k \in \mathcal
C(\intcc01^2)$ et, pour $f \in \mathcal C(\intcc01)$, $Tf(x)
= \int_0^1k(x,y)f(y)\,\dd y$.
(a) Montrer que $T$ envoie la boule unité de $\mathcal
C(\intcc01)$ sur un ensemble équicontinu et uniformément
borné~; conclure que $T$ est un \emph{opérateur compact}~: les
images d'ensembles bornés sont relativement compactes.
(b) En déduire que si $(f_n)$ est bornée, $(Tf_n)$ a une
sous-suite uniformément convergente, et que $T$ ne peut être
une bijection d'inverse \omterm{def:b3:topology:continuity}{continue}. \emph{(L'image de la boule
unité serait un \omterm{def:b3:topology:topology}{voisinage} compact de $0$ dans $\mathcal
C(\intcc01)$~: interdit par le théorème de Riesz de deuxième
année.)}
\end{exercise}

\begin{exercise}[$\star\star$]\label{exo:b3:complete:8}
Prouver ou infirmer, pour $f_n \colon \intcc01 \to \R$
\omterm{def:b3:topology:continuity}{continues}~:
(a) $f_n \downarrow 0$ ponctuellement $\Rightarrow$
uniformément (Dini --- le reprouver)~; (b) même chose sans
monotonie~; (c) même chose avec monotonie mais $f$
discontinue~; (d) même chose avec monotonie, limite \omterm{def:b3:topology:continuity}{continue},
mais sur $\intoo01$.
\end{exercise}

\begin{exercise}[$\star\star$]\label{exo:b3:complete:9}
(a) (Les moments déterminent) Soit $f \in \mathcal
C(\intcc01)$ avec $\int_0^1 f(x)\,x^n\,\dd x = 0$ pour tout $n
\in \N$. Montrer $f = 0$. \emph{(Approximer $f$ uniformément
par des polynômes et calculer $\int f^2$.)}
(b) Montrer que les polynômes pairs sont \omterm{def:b3:topology:interior}{denses} dans $\mathcal
C(\intcc01)$ mais \emph{pas} dans $\mathcal C(\intcc{-1}1)$~;
où l'hypothèse de Stone--Weierstrass échoue-t-elle~
(c) L'algèbre engendrée par $x \mapsto \eu^{\iu x}$ seule
(sans $\eu^{-\iu x}$) est-elle \omterm{def:b3:topology:interior}{dense} dans $\mathcal C(S^1,
\C)$~ \emph{(Considérer $\int_0^{2\pi}f(t)\,\eu^{\iu t}\dd
t$.)}
\end{exercise}

\begin{exercise}[$\star\star\star$]\label{exo:b3:complete:10}
(Bornitude uniforme, version métrique) Soit $X$ un \omterm{def:b3:complete:complete}{espace
métrique complet} et $\mathcal F \subseteq \mathcal C(X, \R)$
une famille \emph{ponctuellement} bornée~:
$\sup_{f \in \mathcal F}\abs{f(x)} < \infty$ pour chaque $x$.
Montrer qu'il existe un \omterm{def:b3:topology:topology}{ouvert} non vide $U \subseteq X$ sur
lequel $\mathcal F$ est \emph{uniformément} bornée~:
$\sup_{f}\sup_U \abs f < \infty$.
\emph{(Considérer $F_n = \{x : \abs{f(x)} \leq n\ \forall
f\}$.)} C'est le moteur de Banach--Steinhaus dans
le~\cref{ch:b3:banach}.
\end{exercise}

\begin{exercise}[$\star\star$]\label{exo:b3:complete:11}
($\mathcal C^1$ a besoin de sa propre norme) Sur $E =
\mathcal C^1(\intcc01, \R)$ considérer $\norm f_{\mathcal C^1}
= \norm f_\infty + \norm{f'}_\infty$.
(a) Montrer que $(E, \norm\cdot_{\mathcal C^1})$ est \omterm{def:b3:complete:complete}{complet}
\emph{(une suite $\mathcal C^1$-Cauchy a $f_n \to f$ et $f_n'
\to g$ uniformément~; identifier $g = f'$ en passant à la
limite dans $f_n(x) = f_n(0) + \int_0^xf_n'$)}.
(b) Montrer que $(E, \norm\cdot_\infty)$ n'est \emph{pas}
\omterm{def:b3:complete:complete}{complet}~: exhiber une limite uniforme de fonctions $\mathcal
C^1$ qui n'est pas dérivable (p.\,ex.\ des approximations
lisses de $\abs{x - \tfrac12}$).
(c) Déduire de (a), (b) et du cercle d'idées de l'application
ouverte --- ou directement --- qu'aucune constante $C$ ne
satisfait $\norm{f'}_\infty \leq C\,\norm f_\infty$ sur $E$~:
exhiber une suite le prouvant. La dérivation est non bornée~;
c'est la falaise derrière
\cref{thm:b3:complete:nowherediff}.
\end{exercise}

\begin{exercise}[$\star\star\star$]\label{exo:b3:complete:12}
(Lemme de Croft) Soit $f\colon\intoo0\infty\to\R$ \omterm{def:b3:topology:continuity}{continue} et
supposons que pour tout $x > 0$, $f(nx) \to 0$ quand l'entier
$n \to \infty$. Montrer que $f(t) \to 0$ quand $t \to
+\infty$. \emph{(Fixer $\varepsilon > 0$~; les ensembles
\[
F_N = \{x > 0 : \abs{f(nx)} \leq \varepsilon\ \ \forall n
\geq N\}
\]
sont fermés et recouvrent $\intoo0\infty$~; Baire dans quelque
intervalle $\intcc ab$ donne $N$ et un sous-intervalle
$\intcc{a'}{b'} \subseteq F_N$~; alors les dilatés
$\bigl[na', nb'\bigr]$, $n \geq N$, recouvrent tout un
\omterm{def:b3:topology:topology}{voisinage} de $+\infty$ une fois $n(b' - a') \geq a'$.)} Où
l'hypothèse «~pour tout $x$~» (pas seulement $x$ rationnel)
est-elle utilisée~
\end{exercise}

\section{Problème~: le théorème d'existence de Peano}

\begin{problem}[{Problème de fin de semaine --- existence de
solutions de $x' = f(t, x)$ sans Lipschitz}]\label{pb:b3:complete:1}
Cauchy--Lipschitz (deuxième année~; reprouvé dans
le~\cref{ch:b3:ode}) exige $f$ lipschitzienne en $x$. Peano
(1890)~: \emph{la continuité de $f$ suffit déjà pour
l'existence} --- mais pas l'unicité. On le prouve avec les
polygones d'Euler et Ascoli. Cadre~: $f \colon R \to \R^d$
\omterm{def:b3:topology:continuity}{continue} sur le rectangle $R = \intcc{t_0 - a}{t_0 + a}\times
\bar B(x_0, b) \subseteq \R\times\R^d$, $M = \sup_R\norm f$,
et
\[
T = \min\bigl(a,\ b/M\bigr)
\qquad (\text{si } M = 0 \text{ le problème est trivial}).
\]

\medskip
\noindent\textbf{Partie I --- Polygones d'Euler.} Pour $n \geq
1$, subdiviser $[t_0, t_0 + T]$ par $t_k = t_0 + kT/n$ ($0
\leq k \leq n$) et définir $\varphi_n$ affine par morceaux~:
$\varphi_n(t_0) = x_0$ et, sur $[t_k, t_{k+1}]$,
\[
\varphi_n(t) = \varphi_n(t_k) + (t -
t_k)\,f\bigl(t_k, \varphi_n(t_k)\bigr).
\]
\begin{enumerate}
  \item Montrer par induction que $\varphi_n$ est bien définie,
        avec $\norm{\varphi_n(t) - x_0} \leq M(t - t_0) \leq b$
        sur $[t_0, t_0 + T]$ --- donc les points d'évaluation
        restent dans $R$. \emph{(C'est là qu'entre $T \leq
        b/M$.)}
  \item Montrer que chaque $\varphi_n$ est $M$-lipschitzienne.
  \item Déduire d'Arzelà--Ascoli
        (\cref{thm:b3:complete:ascoli}) qu'une sous-suite
        $(\varphi_{n_j})$ converge uniformément sur $[t_0, t_0
        + T]$ vers quelque $\varphi$, elle-même
        $M$-lipschitzienne avec $\varphi(t_0) = x_0$.
\end{enumerate}

\medskip
\noindent\textbf{Partie II --- La limite résout l'équation.}
Définir le \emph{défaut} $\Delta_n(t) = \varphi_n'(t) -
f\bigl(t, \varphi_n(t)\bigr)$ hors des nœuds de grille.
\begin{enumerate}[resume]
  \item Montrer que $f$ est uniformément \omterm{def:b3:topology:continuity}{continue} sur $R$, et
        en déduire~: pour tout $\varepsilon > 0$ il existe
        $n_0$ tel que pour $n \geq n_0$ et tout $t$ non nœud,
        $\norm{\Delta_n(t)} \leq \varepsilon$. \emph{(Sur
        $(t_k, t_{k+1})$, $\varphi_n'(t) = f(t_k,
        \varphi_n(t_k))$, et $(t, \varphi_n(t))$ est à
        distance $(1 + M)\,T/n$ de $(t_k,
        \varphi_n(t_k))$.)}
  \item Établir la forme intégrale~: pour tout $t$,
        \[
        \varphi_n(t) = x_0 + \int_{t_0}^{t}
        f\bigl(s, \varphi_n(s)\bigr)\dd s +
        \int_{t_0}^t \Delta_n(s)\,\dd s,
        \]
        l'intégrande du milieu étant \omterm{def:b3:topology:continuity}{continue} par morceaux.
  \item Passer à la limite le long de $(n_j)$~: montrer
        $f(s, \varphi_{n_j}(s)) \to f(s, \varphi(s))$
        \emph{uniformément} (continuité uniforme de $f$
        encore), et conclure
        \[
        \varphi(t) = x_0 + \int_{t_0}^t
        f\bigl(s, \varphi(s)\bigr)\dd s .
        \]
  \item En déduire que $\varphi$ est $\mathcal C^1$ sur $[t_0,
        t_0 + T]$ et résout $x' = f(t, x)$, $x(t_0) = x_0$~;
        étendre la construction à $[t_0 - T, t_0]$ (renversement
        du temps). \emph{C'est le théorème de Peano.}
\end{enumerate}

\medskip
\noindent\textbf{Partie III --- L'unicité échoue vraiment.}
Considérer $x' = 2\sqrt{\abs x}$, $x(0) = 0$, sur $\R$.
\begin{enumerate}[resume]
  \item Vérifier que $f(x) = 2\sqrt{\abs x}$ est \omterm{def:b3:topology:continuity}{continue} mais
        pas lipschitzienne sur aucun \omterm{def:b3:topology:topology}{voisinage} de $0$.
  \item Vérifier que $x \equiv 0$ et, pour tout $c \geq 0$,
        \[
        x_c(t) = \begin{cases} 0 & t \leq c,\\ (t - c)^2 & t
        \geq c,\end{cases}
        \]
        sont toutes solutions passant par $(0,0)$~: un
        continuum de solutions distinctes.
  \item Où l'argument d'itération de Picard (point fixe de
        Banach) casse-t-il pour cette $f$~
\end{enumerate}

\medskip
\noindent\textbf{Partie IV --- Limites de la méthode.}
\begin{enumerate}[resume]
  \item Montrer que le théorème de Peano échoue en dimension
        infinie~: on admet (ou on peut prendre sur foi)
        l'exemple classique de Dieudonné dans l'espace $c_0$
        des suites nulles~; à la place, prouver l'ingrédient
        de dimension finie qui y échoue~: la boule unité
        fermée de $c_0$ (norme sup) n'est pas compacte ---
        exhiber une suite bornée sans sous-suite convergente,
        et expliquer quelle étape de la partie~I casse.
  \item Résumer~: quelles hypothèses donnent l'existence~
        l'existence et l'unicité~? Énoncer précisément les
        deux théorèmes (Peano~; Cauchy--Lipschitz) côte à
        côte.
\end{enumerate}

\medskip
\noindent\textbf{Partie V --- Osgood~: unicité au-delà de
Lipschitz.} Soit $\omega\colon\intoo0\infty\to\intoo0\infty$
\omterm{def:b3:topology:continuity}{continue}, croissante, avec
\[
\int_0^1\frac{\dd r}{\omega(r)} = +\infty
\qquad(\text{divergence en } 0),
\]
et supposons que $f$ satisfait $\norm{f(t, x) - f(t, y)} \leq
\omega\bigl(\norm{x - y}\bigr)$ sur $R$.
\begin{enumerate}[resume]
  \item Vérifier que $\omega(r) = Lr$ qualifie (Lipschitz), que
        $\omega(r) = r\log\frac1r$ (étendue par continuité,
        pour $r$ petit) qualifie bien qu'elle ne soit
        \emph{pas} $O(r)$, et que $\omega(r) = 2\sqrt r$ ne
        qualifie pas. Calculer l'intégrale dans chaque cas.
  \item Soient $x_1, x_2$ solutions de l'équation sur $[t_0,
        t_0 + T]$ avec $x_1(t_0) = x_2(t_0)$, et $\delta(t) =
        \norm{x_1(t) - x_2(t)}$. Montrer, à partir des formes
        intégrales seules, que pour $t_0 \leq s \leq t$~:
        \[
        \delta(t) \leq \delta(s) +
        \int_s^t\omega\bigl(\delta(v)\bigr)\dd v .
        \]
  \item (Théorème d'Osgood) Supposer $\delta(t_1) > 0$ pour
        quelque $t_1$, et soit $\tau = \sup\{t \leq t_1 :
        \delta(t) = 0\}$. Pour $s \in \intoo\tau{t_1}$, poser
        $u(t) = \delta(s) + \int_s^t\omega(\delta(v))\,\dd v$.
        Montrer $\delta \leq u$, $u' = \omega(\delta) \leq
        \omega(u)$, et en déduire
        \[
        \int_{\delta(s)}^{u(t_1)}\frac{\dd r}{\omega(r)}
        \leq t_1 - s \leq t_1 - \tau .
        \]
        Faire $s \downarrow \tau$ et dériver une contradiction
        avec la divergence de l'intégrale. Conclure~: les
        \emph{solutions passant par une condition initiale
        commune coïncident} --- unicité sous la condition
        d'Osgood.
  \item Tirer les conséquences~: l'unicité Cauchy--Lipschitz
        est le cas $\omega(r) = Lr$~; l'équation $x' =
        x\log\frac1{\abs x}$ (étendue par $0$ en $0$) a des
        solutions uniques bien que son second membre ne soit
        pas lipschitzien en $0$~; et pour $x' = 2\sqrt{\abs x}$
        la convergence de $\int_0\frac{\dd r}{2\sqrt r}$ est
        exactement ce qui laisse une solution quitter $0$ en
        temps fini --- apparier la valeur de l'intégrale
        $\int_0^{h^2}\frac{\dd r}{2\sqrt r} = h$ avec le
        comportement d'échappement de $x_c$.
\end{enumerate}

\medskip
\noindent\textbf{Partie VI --- Taux, schémas, entonnoirs.}
\begin{enumerate}[resume]
  \item (Lemme de Grönwall intégral) Soient $e, \eta \geq 0$
        \omterm{def:b3:topology:continuity}{continues} sur $[t_0, t_0+T]$ avec $e(t) \leq
        \int_{t_0}^t\bigl(L\,e(s) + \eta(s)\bigr)\dd s$ pour
        tout $t$. Montrer
        \[
        e(t) \leq \frac{\sup\eta}{L}\bigl(\eu^{L(t - t_0)} -
        1\bigr)
        \]
        \emph{(poser $G(t) = \int_{t_0}^t(Le + \eta)$, noter
        $G' \leq LG + \sup\eta$, et dériver
        $\eu^{-Lt}G(t)$)}.
  \item (Euler converge avec un taux) Supposer maintenant $f$
        $L$-lipschitzienne en $x$ et $L'$-lipschitzienne en
        $t$ sur $R$. Combinant la borne de défaut de la
        question~4 (rendue quantitative~: $\norm{\Delta_n}
        \leq (L' + LM)T/n$) avec la question~17, prouver
        \[
        \sup_{[t_0, t_0+T]}\norm{\varphi_n - \varphi}
        \;\leq\; \frac{(L' + LM)\,T}{n}\cdot
        \frac{\eu^{LT} - 1}{L} = O\Bigl(\frac1n\Bigr),
        \]
        où $\varphi$ est \emph{la} solution~: avec des données
        lipschitziennes toute la suite converge, avec un taux
        explicite --- pas besoin de sous-suites. Pourquoi
        l'unicité promeut-elle la convergence sous-séquentielle
        en convergence pleine même sans ce calcul~
  \item (Le schéma choisit) Pour $x' = 2\sqrt{\abs x}$, $x(0) =
        0$~: montrer que tout polygone d'Euler est
        identiquement zéro, donc le schéma converge vers la
        solution $x \equiv 0$~; mais démarré en $x(0) =
        \varepsilon > 0$ il converge (quand $n \to \infty$,
        puis $\varepsilon \to 0$) vers $t \mapsto t^2$, une
        solution \emph{différente} passant par l'origine. La
        non-unicité resurgit comme sensibilité du schéma
        numérique aux perturbations.
  \item Montrer que l'ensemble $\mathcal S$ de \emph{toutes}
        les solutions de $x' = f(t,x)$, $x(t_0) = x_0$ sur
        $[t_0, t_0+T]$ (valeurs dans $\bar B(x_0, b)$) est non
        vide (partie~II), uniformément $M$-lipschitzien, et
        fermé dans $\bigl(\mathcal C([t_0, t_0+T], \R^d),
        \norm\cdot_\infty\bigr)$~; conclure par Ascoli que
        $\mathcal S$ est \emph{compact}. (Le théorème de
        Kneser ajoute que $\mathcal S$ est \omterm{def:b3:topology:connected}{connexe}~; on ne le
        prouvera pas.)
  \item Vérifier le phénomène de Kneser sur l'exemple~: pour
        $x' = 2\sqrt{\abs x}$, $x(0) = 0$, sur $[0, 1]$,
        montrer $\mathcal S = \{x_c : c \in \intcc0\infty\}$
        avec $x_\infty \equiv 0$ \emph{(pour toute solution,
        poser $c = \sup\{t : x(t) = 0\}$ et intégrer $(\sqrt
        x)' = 1$ sur $\{x > 0\}$)}, que $c \mapsto x_c$ est
        \omterm{def:b3:topology:continuity}{continue} de $\intcc0\infty$ (\omterm{def:b3:topology:locallycompact}{compactification
        d'Alexandroff}, c'est-à-dire $c = \infty$ collé comme
        limite) vers $\mathcal C([0,1])$, et conclure que
        $\mathcal S$ est bien compact et \omterm{def:b3:topology:connected}{connexe} --- un
        entonnoir en forme de segment.
  \item (Ensembles atteignables) Déduire de la question~20 que
        pour chaque $t$ fixé, l'\emph{ensemble atteignable}
        $\mathcal S(t) = \{x(t) : x \in \mathcal S\}$ est
        compact~; le calculer pour l'exemple de la question~21
        et vérifier qu'il est aussi \omterm{def:b3:topology:connected}{connexe}~: $\mathcal S(t)
        = \intcc0{t^2}$ --- tout état intermédiaire est
        atteint par quelque solution.
\end{enumerate}

\medskip
\noindent\textbf{Partie VII --- Compléments~: dépendance,
optimalité, et un schéma calculé à la main.}
\begin{enumerate}[resume]
  \item (Dépendance \omterm{def:b3:topology:continuity}{continue}) Supposer $f$ $L$-lipschitzienne
        en $x$ sur $R$, et soient $x, y$ deux solutions de
        valeurs initiales $x_0, y_0$ en $t_0$. En adaptant la
        preuve de la question~17 à l'inégalité $e(t) \leq
        \norm{x_0 - y_0} + \int_{t_0}^tL\,e(s)\dd s$, montrer
        \[
        \norm{x(t) - y(t)} \leq \norm{x_0 - y_0}\,
        \eu^{L(t - t_0)},
        \]
        et vérifier sur $x' = Lx$ que la borne est atteinte~:
        Grönwall est optimal. En déduire à nouveau l'unicité
        ($x_0 = y_0$), et que le flot $x_0 \mapsto x(t; x_0)$
        est lipschitzien, de constante $\eu^{LT}$, là où il est
        défini.
  \item (Osgood est optimal) Réciproquement, soit $\omega$
        \omterm{def:b3:topology:continuity}{continue}, croissante, positive sur $\intoo0\infty$,
        avec $\omega(r) \to 0$ quand $r \to 0^+$ mais
        \[
        \int_0^1\frac{\dd r}{\omega(r)} < +\infty,
        \]
        et étendre $f(x) = \omega(\abs x)$ par $f(0) = 0$.
        Montrer que $\Omega(x) = \int_0^x\frac{\dd
        r}{\omega(r)}$ est une bijection croissante de
        $\intoc0{1}$ sur $\intoc0{\Omega(1)}$, que son inverse
        $g$ résout $g' = \omega(g)$ avec $g(0^+) = 0$, et que
        $g$, étendue par $0$ pour $t \leq 0$, est une solution
        $\mathcal C^1$ de $x' = f(x)$ passant par $(0, 0)$
        distincte de $x \equiv 0$ \emph{(pour $g'(0) = 0$,
        borner $\frac{g(t)}t$ par $\omega(g(t))$)}. Conclure~:
        l'hypothèse de divergence de la question~15 n'est pas
        une commodité mais la frontière exacte de l'unicité~;
        retrouver la partie~III depuis $\omega(r) = 2\sqrt r$,
        $\Omega(x) = \sqrt x$.
  \item (Euler calculé à la main) Pour $x' = x$, $x(0) = 1$ sur
        $[0, 1]$~: montrer que le polygone d'Euler satisfait
        $\varphi_n(1) = \bigl(1 + \frac1n\bigr)^n$. Prouver le
        développement
        \[
        \Bigl(1 + \frac1n\Bigr)^{n}
        = \eu\Bigl(1 - \frac1{2n} + O\Bigl(\frac1{n^2}\Bigr)
        \Bigr),
        \]
        donc l'erreur en $t = 1$ est $\frac{\eu}{2n} +
        O(n^{-2})$~: le taux $O(\frac1n)$ de la question~18,
        avec la constante exacte. Vérifier numériquement pour
        $n = 10$~: $1.1^{10} = 2.59374$ contre $\eu \approx
        2.71828$, une erreur $0.12454$ à comparer avec
        $\frac{\eu}{20} \approx 0.13591$.
\end{enumerate}
\end{problem}
